\section{C1. Normalized observations, complete cylinder-dual spaces, and an exact adaptive diagnostic experiment}
\label{sec:r33-c1}

The original objective is inference from partially observed deterministic
mechanics.  Three logically different assertions must be separated: a valid
probability channel, finite-horizon differentiability of its filter, and
long-time statistical identification.  The first two do not imply a strict
filter contraction.  We prove them on a specified infinite product state and
provide an exactly solvable adaptive diagnostic experiment for the third.
This example is not asserted to identify arbitrary mechanical parameters.

\subsection{The probability of selecting a stratum}

Let $Y=\bigsqcup_{s\ge1}Y_s$, let $m_s$ be a sigma-finite measure, and let
$q_s>0$ be an $m_s$-density of total mass one.  Fix $\alpha_s>0$ with
$\sum_s\alpha_s=1$ and put $\nu|_{Y_s}=\alpha_sq_sm_s$.
For each hidden state, parameter, and action, specify measurable numbers
$p_s(\theta,a,x)\ge0$ of sum one and conditional densities
$\rho_s(\theta,a,x;\cdot)$ of mass one relative to $m_s$.  Then
\begin{equation}\label{eq:r33-c1-density}
 g_\theta^a(x,y)=
 \frac{p_s(\theta,a,x)\rho_s(\theta,a,x;y)}{\alpha_sq_s(y)},
 \qquad y\in Y_s.
\end{equation}

\begin{proposition}[Whole-channel normalization]
\label{prop:r33-c1-channel}
The measure $g_\theta^a(x,y)\nu(dy)$ is a probability kernel.  If the
parameter derivatives of $p_s\rho_s$ through order $r$ are continuous in
$L^1(m_s)$ and their $L^1$ norms are summable uniformly on compact parameter
sets, the channel is $C^r$ in $L^1(\nu)$.
\end{proposition}
\begin{proof}
Multiplication by $\nu$ cancels $\alpha_sq_s$ on stratum $s$ and its integral
is $p_s$.  Tonelli gives total mass $\sum_sp_s=1$.  The same cancellation
identifies the $L^1(\nu)$ derivative norm with the sum of the $L^1(m_s)$
norms, so the Banach-valued Weierstrass theorem applies.  This does not claim
quadratic-mean differentiability at a parameter where a previously absent
stratum acquires positive probability; square-root differentiability requires
its own condition.
\end{proof}

\subsection{A complete cylinder-dual space with explicit bonding maps}

For this construction take the genuinely infinite-dimensional hidden space
$X=(\mathbb T^d)^{\mathbb N}$ with its product topology.  It is compact
metrizable.  Let $s>d/2$ and, for $j\ge0$, set
\[
 H_{n,j}=\widehat\bigotimes_{i=1}^n H^{-(s+j)}(\mathbb T^d).
\]
All tensor products are Hilbert tensor products with Fourier weights.
The map $b_{n+1,n}:H_{n+1,j}\to H_{n,j}$ is contraction with the constant
function $1$ in the last coordinate, whose positive Sobolev norm is one.
Define
\[
 \mathscr E_j=
 \{(T_n)_{n\ge1}:b_{n+1,n}T_{n+1}=T_n\},
 \qquad p_{n,j}(T)=\|T_n\|_{H_{n,j}}.
\]

\begin{proposition}[Complete projective dual and probability embedding]
\label{prop:r33-c1-projective}
The space $\mathscr E_j$ is a complete metrizable locally convex space, with
metric $\sum_n2^{-n}(1\wedge p_{n,j}(T-S))$.  Finite signed measures on $X$
embed by their finite-coordinate marginals.  The embedding of probability
measures is continuous from weak convergence into every seminorm $p_{n,j}$.
\end{proposition}
\begin{proof}
The countable product of the Hilbert spaces is complete and the compatibility
equations are kernels of continuous linear maps.  Thus their intersection is
closed and complete.  For one coordinate,
$\sum_{k\in\mathbb Z^d}(1+|k|^2)^{-s-j}<\infty$ makes $x\mapsto\delta_x$
a continuous map into $H^{-(s+j)}$.  Finite tensor products give the same
property for $x_{1:n}\mapsto\delta_{x_{1:n}}$, with compact image.  Approximate
this Hilbert-valued continuous map uniformly by finite-rank continuous maps.
Weak convergence of marginals then implies norm convergence of their Bochner
integrals, which are precisely the marginal distributions in $H_{n,j}$.
Bonding compatibility follows from marginalization.
\end{proof}

Suppose $\Phi_\theta:X\to X$ is compatible and its first $n$ output
coordinates depend on finitely many input coordinates, with $C^{r+1}$
parameter derivatives bounded on compact parameter sets.  No inverse map is
required.  Spatial differentiability is needed only for the finite output
coordinates tested below.

\begin{theorem}[Strong parameter derivatives on the common cylinder level]
\label{thm:r33-c1-push}
For a fixed finite signed measure $\mu$ on $X$, the map
$\theta\mapsto(\Phi_\theta)_\#\mu$ is $C^r$ in $\mathscr E_r$.
For a smooth cylindrical test $\varphi$,
\[
 \langle\partial_\theta^j(\Phi_\theta)_\#\mu,\varphi\rangle
 =\int_X\partial_\theta^j[\varphi(\Phi_\theta(x))]\,\mu(dx),
 \qquad j\le r.
\]
\end{theorem}
\begin{proof}
For a finite coordinate, the Fourier series for $\partial^\ell\delta_x$ is
square summable in $H^{-(s+r)}$ for $\ell\le r$; dominated convergence of
that series makes these derivatives norm continuous.  The finite tensor
product has norm-continuous derivatives of total order at most $r$.
Composition with the first $n$ output coordinates of $\Phi_\theta$ therefore
has a Taylor remainder tending to zero uniformly in $x$ on each compact
parameter set.  Integrating that remainder against $|\mu|$ proves strong,
not merely weak, differentiability in $H_{n,r}$.  The identities are compatible
under every $b_{n+1,n}$, so convergence holds in the stated projective topology.
The assertion is for measures and their resulting jets; it does not assert
that an arbitrary nonlinear push-forward is bounded on every distributional
dual space.
\end{proof}

For a fixed observation record of length $n$, write the unnormalized posterior
as one integral over the initial hidden state:
\[
 \langle N_{n,\theta},\varphi\rangle
 =\int\varphi(X_n^\theta(x))
   \prod_{k=1}^n g_\theta^{a_k}(X_k^\theta(x),y_k)\,\mu(dx).
\]
If its likelihood factors and output coordinate maps have the preceding
finite-coordinate smoothness and bounded derivative envelopes, the same
Taylor argument differentiates $N_{n,\theta}$ in each $H_{m,r}$.
Where $r_{n,\theta}=N_{n,\theta}(1)>0$, normalization gives
\[
 \partial_\theta\Pi_{n,\theta}
 =r_{n,\theta}^{-1}\partial_\theta N_{n,\theta}
  -r_{n,\theta}^{-2}N_{n,\theta}\partial_\theta r_{n,\theta}.
\]
These are finite-horizon estimates; their constants can grow with $n$.
A positive likelihood and a deterministic permutation do not imply a strict
Hilbert-projective contraction: multiplication cancels in posterior density
ratios.  No such contraction is used here.

\subsection{Nonoverlapping diagnostic scheduling}

A policy chooses an action using the observed past and a fresh external
randomizer before the next noise is drawn.  Its conditional action law is a
fixed measurable function of that past and does not explicitly depend on the
unknown parameter.  To schedule length-$m_0$ diagnostic words, partition time
into disjoint blocks.  At a block start choose one word $j$ with probability
at least $\rho_j$, where $\sum_j\rho_j\le1$, and execute that word until the
block ends.  Distinct incompatible words are not required to occur in the
same block.  This defines a nonempty causal policy class.

\subsection{An exact, nonempty adaptive statistical example}

Let the hidden state evolve by any parameter-free deterministic map of $X$.
The parameter $\theta\in\mathbb R^d$ is a diagnostic sensor offset.  At time
$k$ the policy selects a diagnostic indicator $A_k\in\{0,1\}$ with
\[
 \mathbb P(A_k=1\mid\mathcal F_{k-1})\ge\rho>0.
\]
When $A_k=1$, observe $Y_k=\theta+\xi_k$, with independent
$\xi_k\sim N(0,\Sigma)$, $\Sigma>0$, independent also of the policy
randomizer and hidden state.  Nondiagnostic outputs have a parameter-free
law and no dependence on these diagnostic noises.  Put
$T_n=\sum_{k\le n}A_k$ and
$S_n=\sum_{k\le n}A_k(Y_k-\theta_0)$.
Take a prior $N(m_0,V_0)$, $V_0>0$.

\begin{theorem}[Exact adaptive likelihood and posterior]
\label{thm:r33-c1-exact}
For every such policy,
\[
 \log\frac{dP_{\theta_0+h/\sqrt n}}{dP_{\theta_0}}
 =h^T\Sigma^{-1}S_n/\sqrt n
   -\frac{T_n}{2n}h^T\Sigma^{-1}h
\]
exactly.  Conditional on all observations and actions, the posterior is
$N(m_n,V_n)$, where
\[
 V_n=(V_0^{-1}+T_n\Sigma^{-1})^{-1},\qquad
 m_n=V_n\left(V_0^{-1}m_0+
            \Sigma^{-1}\sum_{k\le n}A_kY_k\right).
\]
Furthermore $\mathbb P(T_n<\rho n/2)\le e^{-c_\rho n}$ and, with
$\bar Y_n=T_n^{-1}\sum A_kY_k$ on $T_n>0$,
\[
 d_{\rm TV}\left(
 \mathcal L(\sqrt{T_n}(\theta-\bar Y_n)\mid\mathcal F_n),
 N(0,\Sigma)\right)\longrightarrow0
\]
in $P_{\theta_0}$-probability uniformly over the policy class.
\end{theorem}
\begin{proof}
The common policy likelihood factors cancel in the chronological joint law
of actions and outputs.  The remaining Gaussian factors give the displayed
quadratic log likelihood.  Completing its square with the Gaussian prior
gives $m_n,V_n$, even though the diagnostic choices are adaptive.
The differences $A_k-E(A_k\mid\mathcal F_{k-1})$ are bounded martingale
differences.  Azuma's inequality proves the bound on $T_n$.
The score sum has $E\|S_n\|^2\le n\operatorname{tr}\Sigma$, uniformly in
the policy, so $\bar Y_n-\theta_0=O_P(n^{-1/2})$ on the high-probability
event $T_n\ge\rho n/2$.

The rescaled posterior covariance is $T_nV_n\to\Sigma$ uniformly on this
event.  Its mean is
$\sqrt{T_n}(m_n-\bar Y_n)
 =\sqrt{T_n}V_nV_0^{-1}(m_0-\bar Y_n)=O_P(n^{-1/2})$.
The explicit relative entropy between two nondegenerate Gaussian laws tends
to zero; Pinsker's inequality proves the total-variation conclusion.
No hidden filter contraction, Riccati identity, or unproved annular testing
argument is needed for this example.
\end{proof}

\subsection{Mechanical identification remains a separate theorem}

The exact diagnostic experiment proves that normalized adaptive inference is
nonempty on an infinite-dimensional deterministic hidden state.  It does not
prove that a billiard or hard-sphere parameter can be identified from an
arbitrary chosen signal.  For that application one must establish induced-law
identification, recursive derivative bounds uniform in time, and posterior
testing for the actual mechanical channel.  A pointwise Fisher inequality is
not substituted for any of these statements.
